\documentclass[12pt]{article}
\usepackage[margin=1in]{geometry}
\usepackage{amsmath, amssymb}
\usepackage{enumitem, array}

\usepackage{boxedminipage}
\usepackage{fancyhdr, actuarialangle}
\usepackage{ragged2e, comment}


\newenvironment{problem}[1][]{%
  \bigskip% put space before problem statement box %
  \noindent\begin{boxedminipage}{\columnwidth}%
  \def\@tempa{#1}%
  \ifx\@tempa\empty\else%
    \textbf{#1}\hspace{0.5em}%
  \fi%
}{%
  \end{boxedminipage}%
}


\begin{document}

\begin{center}
    \Large \textbf{MATH 2790 Test 3 Fall 2025} \\[1em]
\end{center}

\vspace{.25cm}

\noindent Name: \underline{\hspace{8cm}} \hspace{.25cm} Auburn ID: \underline{\hspace{4cm}}



\section*{Instructions}
\begin{enumerate}[label=\arabic*.]
\item Print your full name and your banner ID on your exam. Then finish reading these instructions, and sign the bottom of this page.
\item All electronic equipment (cell phones, Apple Watch, AirPod, etc) are \textbf{NOT ALLOWED} during the exam time.
\item You are allowed to use up to 2 calculators from the following list: BA-35, BA II Plus, BA II Plus Professional, TI-30Xa, TI-30X II (IIS solar or IIB battery), TI-30XS MultiView (or XB battery).
\item No talking is allowed during the exam.
\item Unless otherwise indicated, \textbf{SHOW ALL YOUR WORK}. If no work is shown, no partial credit can be awarded. Your work and answer need to be accurate and relevant to receive points.
\item You will be given exactly 50 minutes for this exam.
\item Any student not following the above instructions nor behaving according to the above instructions during the exam may have their exam confiscated and points deducted.
\end{enumerate}

I have read and fully understand all of the above instructions. \\[2em]
Signature: \hrulefill


\vspace{.75cm}
Do not write in the table below

% scoring table
\begin{center}
\begin{tabular}{|*{3}{>{\centering\arraybackslash}p{.1\linewidth}|}}
    \hline
    \textbf{Problem} & \textbf{Points} & \textbf{Score} \\
    \hline
    1 & 15 &  \\
    \hline
    2 & 15 &  \\
    \hline
    3 & 15 &  \\
    \hline
    4 & 15 &  \\
    \hline
    5 & 16 &  \\
    \hline
    6 & 12 &  \\
    \hline
    7 & 12 &  \\
    \hline \hline
    \textbf{Total} & 95+5 &  \\
    \hline
\end{tabular}
\end{center}

\newpage

%%%%% PROBLEM 1 %%%%%
\begin{problem}[1. (16 points; 4 each)]
Marcus plans to buy an apartment unit in Atlanta. He is able to make a down
payment of \$60{,}000 and plans to apply for a 25-year mortgage. The annual effective
interest rate is 6.8\% and he can afford a monthly payment of \$2{,}000. Find the value in
each following part, \textbf{rounding your \underline{final} answers to 2 decimal places}.

\begin{enumerate}
    \item[(a)] Find the price of the most expensive condo Marcus can afford.
    \item[(b)] Find the amount of interest paid in the 45th payment.
    \item[(c)] Find the outstanding loan balance after 8 years.
    \item[(d)] Find the total amount of interest paid over the first 8 years.
\end{enumerate}
\end{problem}


% %%%%% PROBLEM 1 %%%%%
% \begin{problem}[1. (xxx points)]
% Lisa plans to buy a single family house in Auburn. She is able to make a down
% payment of \$50{,}000 and plans to apply for a 30-year mortgage. The annual effective
% interest rate is 7.5\% and she can afford a monthly payment of \$1{,}500. Find the value in
% each following part, rounding answers to 4 decimal places.

% \begin{enumerate}
%     \item[(a)] Find the price of the most expensive house Lisa can afford.
%     \item[(b)] Find the amount of interest paid in the 28th payment.
%     \item[(c)] Find the outstanding loan balance after 10 years.
%     \item[(d)] Find the total amount of interest paid over the first 10 years.
% \end{enumerate}
% \end{problem}


\newpage


%%%%% PROBLEM 2 %%%%%

\begin{problem}[2. (16 points)]
An 18-year bond is purchased at a discount. The bond pays semiannual coupons.
The amount for accumulation of discount in the 15th coupon is 90.01. The amount for 
accumulation of discount in the 31st coupon is 176.52.

\begin{enumerate}
    \item[(a)] Calculate the discount of this bond. Round to 4 decimal places.
    \item[(b)] Calculate the total amount for accumulation of discount in the first 10 years. Round to 2 decimal places.
\end{enumerate}
\end{problem}


% %%%%% PROBLEM 2 %%%%%

% \begin{problem}[2. (18 points)]
% A 40-year bond is purchased at a discount. The bond pays annual coupons. The
% amount for accumulation of discount in the 15th coupon is 194.82. The amount for
% accumulation of discount in the 20th coupon is 306.69.

% \begin{enumerate}
%     \item[(a)] Calculate the discount of this bond. Round to 4 decimal places.
%     \item[(b)] Calculate the total amount for accumulation of discount in the first 20 years.
% \end{enumerate}
% \end{problem}


\newpage

%%%%% PROBLEM 3 %%%%%
\begin{problem}[3. (12 points)]
Consider two 30-year bonds with the same purchase price. Each has a
nominal coupon rate of 5\% paid semiannually and a face value of \$1{,}000.

\medskip

The first bond has a nominal yield rate of 6\% compounded semiannually, and a
redemption value of \$1{,}200.

\medskip

The second bond has an \textit{effective} yield rate of $j$ compounded semiannually, and a
redemption value of \$900.

\medskip

Find the purchase price and calculate $j$. Round your answer for $j$ to 4 decimal places. If you use a financial calculator please write exactly what each TVM variable is.
\end{problem}


\newpage

%%%%% PROBLEM 4 %%%%%
\begin{problem}[4. (15 points)]
Taylor purchases a \$6{,}000 5\% ten-year par-value bond with annual coupons. If held to
maturity, her yield rate is 4.5\%. The bond is callable at the end of three years for \$6{,}200
and at the end of five, six, and seven years for \$6{,}100.

\begin{enumerate}
    \item[(a)] Find the price of the bond. Round to 4 decimal places. 
    \item[(b)] State the fewest number of call dates that may give the minimal yield rate \textit{without
    calculating} the yield rate and explain your reasoning.
    \item[(c)] Based on your statement in part (b), calculate Sam’s minimal yield
    rate for the period he holds the bond. Round to 4 decimal places.
\end{enumerate}
\end{problem}


\newpage

%%%%% PROBLEM 5 %%%%%

\begin{problem}[5. (12 points)]
A loan of \$95{,}104 is to be repaid with payments at the end of each year for twenty
years. The first payment is $X$. Each subsequent payment is 1\% greater than the
preceding payment. The loan payments are based on an effective interest rate of 4\%.

\begin{enumerate}
    \item[(a)] Calculate the loan balance immediately after the 4th payment. Round
    to 4 decimal places.
    \item[(b)] Calculate the interest due in the 5th payment. Round to 2 decimal
    places.
\end{enumerate}
\end{problem}



\newpage

%%%%% PROBLEM 6 %%%%%
\begin{problem}[6. (12 points)]
Brian pays off a loan with level annual payments at the end of each year. The annual
effective interest rate is 5\%. The outstanding balance at the end of the sixth year is
\$36{,}370.33120 and at the end of the seventh year is \$36{,}188.84776.

\begin{enumerate}
    \item[(a)] Calculate the amount of the level payment. Round to the nearest integer. 
    \textit{Hint: the balances given are just one year apart, is there a formula to relate $B(t)$ and $B(t+1)$?}
    \item[(b)] Calculate the total loan amount. Round to 2 decimal places.
\end{enumerate}
\end{problem}


\newpage


%%%%% PROBLEM 7 %%%%%
\begin{problem}[7. (12 points)]
A loan of $\$12,000$ is taken out at an annual effective interest rate of $7\%$. Connor makes annual payments that are three times the amount of interest accrued in the year. He continues this process until he can make one final payment of at most $\$800$. Find the time and the amount of Connor's final payment. Round to 2 decimal places.
\end{problem}

\newpage
Extra page for scratch Work



\end{document}
